% Chapter 7.3 Test Management Process and 7.4 Bugs Report
\section{Test Management Process}
\label{sec:test_management}

\subsection{Design the Test Strategy}
\begin{itemize}[nosep]
  \item \textbf{Use-case based tests.} Every test case belongs to one of the twelve use cases, so the unit test tables follow the same names and order as the use cases.
  \item \textbf{Tests with the feature.} Tests are written in the same sprint as the feature they check.
  \item \textbf{Isolation.} No test uses the live database; each run starts from fresh test data.
  \item \textbf{Defects.} Each defect is recorded, fixed, and the test that found it stays in the suite.
  \item \textbf{Entry and exit criteria.} Entry: the code builds and the test cases are written. Exit: no failed test and no open defect.
\end{itemize}

\subsection{Test Objectives}
\begin{itemize}[nosep]
  \item \textbf{Functional checks:} each requirement FR1--FR12 behaves as specified, including the reason shown when an action is refused.
  \item \textbf{Access control:} each role, exam owner and assigned invigilator sees only their own data.
  \item \textbf{Exam integrity:} the server enforces the start time, the deadline, the seat and a single sitting, and gives a receipt on submission.
  \item \textbf{Resilience:} answers and monitoring events survive network loss, a client restart and a crash.
\end{itemize}

\subsection{Test Criteria}
For grading accuracy and for each non-functional requirement, Table~\ref{tab:test_criteria} gives a measurable criterion. The last column says which test cases in Section~\ref{sec:test_levels} check it, and a criterion that was not measured says so.

{\small\setlength{\tabcolsep}{4pt}\setlength{\doublerulesep}{1pt}
\begin{longtable}{@{}P{2.8cm}P{9.6cm}P{2.9cm}@{}}
\caption{Test Criteria}\label{tab:test_criteria}\\ \hline\hline
\textbf{Test Scenario} & \textbf{Test Criteria Description} & \textbf{Verified by}\\ \hline
\endfirsthead
\multicolumn{3}{c}{\footnotesize\tablename~\thetable{} -- continued from previous page}\\ \hline\hline
\textbf{Test Scenario} & \textbf{Test Criteria Description} & \textbf{Verified by}\\ \hline
\endhead
\hline \multicolumn{3}{r}{\footnotesize Continued on next page}\\
\endfoot
\hline\hline
\endlastfoot
Grading accuracy (FR11) & An MCQ answer gets full or zero marks. A coding answer is marked by test-case weight (weights 1, 1 and 2 with the two weight-1 cases passed give 50\%). A renamed copy scores at least 90\% similarity, a commented copy at least 95\% and a different algorithm below 30\%. & UT-51, UT-52, UT-54\\
NFR1 Security & A wrong password and a second PC are refused; only a super administrator creates administrators; a deactivated or demoted user is signed out at once; a session cannot be renewed after logout. & UT-02, UT-04, UT-34, UT-35, IT-01, ST-01\\
NFR2 Reliability & Offline events keep their order and survive a restart; offline answer saves are replayed in order; a changed backup file is rejected; Task Manager is restored after a crash. & UT-21, UT-27, UT-28, UT-19\\
NFR3 Usability & A failed login shows an error and stays on the login page; every refused start gives a reason the user can read. No formal user study was carried out. & UT-02, IT-03, IT-04, IT-05, ST-04\\
NFR4 Performance & Readiness fails a PC whose server round trip is over 2~s or whose clock differs by more than 60~s; autosave runs 10~s after an edit; a submission returns in under 1~s while grading continues. The 40-client load target was not measured. & UT-06, UT-07, UT-26, IT-11\\
NFR5 Maintainability & Every use case has automated tests that run on one PC, and a full test run takes under 20 minutes (17 minutes in the final run). & Section~\ref{sec:test_levels}\\
NFR6 Data Integrity & An older save cannot replace a newer one; the client and the server compute the same receipt; a submitted exam cannot be started again; a changed or deleted audit row is found at that row. & UT-30, UT-57, UT-58, IT-07, IT-08, IT-12\\
NFR7 Compatibility & The run used Windows~11 Pro and the Chromium browser; Windows~10 and Edge were not tested. & All test cases\\
NFR8 Connectivity & Warnings, unlocks and extra time reach the student client in the heartbeat reply, and the web panel shows them live. HTTPS was not tested. & IT-10, ST-07, ST-08\\
NFR9 Privacy & Screenshots can be opened only by the exam's staff; students and other teachers are refused. & IT-09\\
NFR10 Backup and Recovery & After a client crash the session is resumed and the newest copy of each answer is kept. & UT-29\\
\end{longtable}}

\subsection{Resource Planning}
Each member wrote the tests for the modules they built, the testing owner ran the full test run, and the supervisor reviewed the results at each sprint review (Section~\ref{sec:managerial}). Table~\ref{tab:test_tools} lists the tools used for testing.

\begin{table}[H]
\centering
\caption{Testing Tools}
\label{tab:test_tools}
\vspace{4pt}
{\small\setlength{\doublerulesep}{1pt}
\begin{tabular}{@{}P{4.6cm}P{10.8cm}@{}}
\hline\hline
\textbf{Tool} & \textbf{Use} \\ \hline
xUnit 2.9 & Unit and integration tests of the server and the desktop client \\
Playwright 1.60 & Web request tests and browser tests of the web panel \\
Docker and MinGW & Compiling and running real C++ answers in the grading tests \\
Git & Version control of the code and the tests \\
\hline\hline
\end{tabular}}
\end{table}

\subsection{Plan Test Environment}
\begin{itemize}[nosep]
  \item \textbf{Hardware:} one development PC running Windows~11 Pro.
  \item \textbf{Server:} the SmartExam server with a fresh in-memory test database; it never connects to the live database.
  \item \textbf{Browser:} Chromium without a visible window, one test at a time.
  \item \textbf{Code runner:} Docker and MinGW, both present in the final run.
\end{itemize}

\subsection{Schedule and Estimation}
Tests were written in the same sprint as the module they check (Chapter~\ref{chap:spp}). Table~\ref{tab:test_schedule} maps the test work to the sprints of Table~\ref{tab:timeframe}; Sprint~8 was kept for system tests, the full test run and defect fixing.

\begin{table}[H]
\centering
\caption{Test Schedule}
\label{tab:test_schedule}
\vspace{4pt}
{\small\setlength{\doublerulesep}{1pt}
\begin{tabular}{@{}P{3.4cm}P{12.0cm}@{}}
\hline\hline
\textbf{Sprint (Weeks)} & \textbf{Testing activity} \\ \hline
Sprint 2 (5--8) & UC-01, UC-07: login, device binding and roles \\
Sprint 3 (9--12) & UC-08, UC-09: exam creation, eligibility and seats \\
Sprint 4 (13--16) & UC-03: student dashboard and exam start \\
Sprint 5 (17--20) & UC-05, UC-10: monitoring data and live supervision \\
Sprint 6 (21--24) & UC-06, UC-11: autosave, submission, grading and plagiarism \\
Sprint 7 (25--28) & UC-12: audit log and reports \\
Sprint 8 (29--32) & UC-02, UC-04: readiness and lockdown; system tests, full test run and defect fixing \\
\hline\hline
\end{tabular}}
\end{table}

\section{Bugs Report}\label{sec:bugs}

In total 49 defects (D-01 to D-49) were recorded, and all of them are fixed. Table~\ref{tab:bugs} lists the 21 with the largest effect on exam integrity, data access and reliability, with the use case it affects.

{\small\setlength{\tabcolsep}{4pt}\setlength{\doublerulesep}{1pt}
\begin{longtable}{@{}P{1.5cm}P{1.8cm}P{8.3cm}P{1.7cm}>{\centering\arraybackslash}p{1.3cm}@{}}
\caption{Bugs Report}\label{tab:bugs}\\ \hline\hline
\textbf{Bug ID} & \textbf{Use Case} & \textbf{Description} & \textbf{Severity} & \textbf{Status}\\ \hline
\endfirsthead
\multicolumn{5}{c}{\footnotesize\tablename~\thetable{} -- continued from previous page}\\ \hline\hline
\textbf{Bug ID} & \textbf{Use Case} & \textbf{Description} & \textbf{Severity} & \textbf{Status}\\ \hline
\endhead
\hline \multicolumn{5}{r}{\footnotesize Continued on next page}\\
\endfoot
\hline\hline
\multicolumn{5}{@{}p{\linewidth}@{}}{\footnotesize Severity: Critical = a leak of exam content or passwords, or a way around the desktop client; High = a broken access rule or a loss of exam time or data.}\\
\endlastfoot
D-01 & UC-08 & Exam details showed the students' password hashes. & Critical & Fixed\\
D-02 & UC-03 & Starting from the desktop client did not mark the exam as active. & High & Fixed\\
D-05 & UC-10 & Invigilator warnings never reached the student client. & High & Fixed\\
D-08 & UC-03 & A session could be started before the exam start time. & High & Fixed\\
D-09 & UC-09 & Each login added its PC to a default lab. & High & Fixed\\
D-10 & UC-09 & Teachers could change other sections' eligibility. & High & Fixed\\
D-12 & UC-01 & Logging out of the web panel left the session valid on the server. & High & Fixed\\
D-14 & UC-01 & A password change left other sessions valid. & High & Fixed\\
D-15 & UC-04 & A failed clean-up step left Task Manager disabled. & High & Fixed\\
D-18 & UC-06 & A session could not be resumed after a crash. & High & Fixed\\
D-25 & UC-05 & Clients could send events that only the server may create. & High & Fixed\\
D-27 & UC-05 & A webcam error closed the exam app. & High & Fixed\\
D-31 & UC-10 & Any teacher could edit or end any exam. & Critical & Fixed\\
D-34 & UC-11 & Any teacher could change the marks of any exam. & Critical & Fixed\\
D-39 & UC-10 & Teachers could open, report on and notify another section's exam. & High & Fixed\\
D-40 & UC-01 & Resetting a device binding left the old PC signed in. & High & Fixed\\
D-41 & UC-10 & A deactivated user kept live-monitor rights on an open connection. & High & Fixed\\
D-42 & UC-03 & Students could read the questions days before the exam. & Critical & Fixed\\
D-43 & UC-06 & A student could sit a submitted exam again. & High & Fixed\\
D-44 & UC-09 & Teachers could add other teachers' students to their exams. & High & Fixed\\
D-45 & UC-09 & A web-panel login could act in the lab PC's exam session. & Critical & Fixed\\
\end{longtable}}

In the final run on 26 October 2026, all 920 executed automated tests passed and none failed. Checks that need real lab PCs (the locked exam window, keyboard and USB events, a real webcam, a pulled network cable and the 40-client load test) were not part of this run and remain on the lab checklist.
